% Options for packages loaded elsewhere
% Options for packages loaded elsewhere
\PassOptionsToPackage{unicode}{hyperref}
\PassOptionsToPackage{hyphens}{url}
\PassOptionsToPackage{dvipsnames,svgnames,x11names}{xcolor}
%
\documentclass[
]{article}
\usepackage{xcolor}
\usepackage[margin=1in]{geometry}
\usepackage{amsmath,amssymb}
\setcounter{secnumdepth}{5}
\usepackage{iftex}
\ifPDFTeX
  \usepackage[T1]{fontenc}
  \usepackage[utf8]{inputenc}
  \usepackage{textcomp} % provide euro and other symbols
\else % if luatex or xetex
  \usepackage{unicode-math} % this also loads fontspec
  \defaultfontfeatures{Scale=MatchLowercase}
  \defaultfontfeatures[\rmfamily]{Ligatures=TeX,Scale=1}
\fi
\usepackage{lmodern}
\ifPDFTeX\else
  % xetex/luatex font selection
\fi
% Use upquote if available, for straight quotes in verbatim environments
\IfFileExists{upquote.sty}{\usepackage{upquote}}{}
\IfFileExists{microtype.sty}{% use microtype if available
  \usepackage[]{microtype}
  \UseMicrotypeSet[protrusion]{basicmath} % disable protrusion for tt fonts
}{}
\makeatletter
\@ifundefined{KOMAClassName}{% if non-KOMA class
  \IfFileExists{parskip.sty}{%
    \usepackage{parskip}
  }{% else
    \setlength{\parindent}{0pt}
    \setlength{\parskip}{6pt plus 2pt minus 1pt}}
}{% if KOMA class
  \KOMAoptions{parskip=half}}
\makeatother
% Make \paragraph and \subparagraph free-standing
\makeatletter
\ifx\paragraph\undefined\else
  \let\oldparagraph\paragraph
  \renewcommand{\paragraph}{
    \@ifstar
      \xxxParagraphStar
      \xxxParagraphNoStar
  }
  \newcommand{\xxxParagraphStar}[1]{\oldparagraph*{#1}\mbox{}}
  \newcommand{\xxxParagraphNoStar}[1]{\oldparagraph{#1}\mbox{}}
\fi
\ifx\subparagraph\undefined\else
  \let\oldsubparagraph\subparagraph
  \renewcommand{\subparagraph}{
    \@ifstar
      \xxxSubParagraphStar
      \xxxSubParagraphNoStar
  }
  \newcommand{\xxxSubParagraphStar}[1]{\oldsubparagraph*{#1}\mbox{}}
  \newcommand{\xxxSubParagraphNoStar}[1]{\oldsubparagraph{#1}\mbox{}}
\fi
\makeatother


\usepackage{longtable,booktabs,array}
\usepackage{calc} % for calculating minipage widths
% Correct order of tables after \paragraph or \subparagraph
\usepackage{etoolbox}
\makeatletter
\patchcmd\longtable{\par}{\if@noskipsec\mbox{}\fi\par}{}{}
\makeatother
% Allow footnotes in longtable head/foot
\IfFileExists{footnotehyper.sty}{\usepackage{footnotehyper}}{\usepackage{footnote}}
\makesavenoteenv{longtable}
\usepackage{graphicx}
\makeatletter
\newsavebox\pandoc@box
\newcommand*\pandocbounded[1]{% scales image to fit in text height/width
  \sbox\pandoc@box{#1}%
  \Gscale@div\@tempa{\textheight}{\dimexpr\ht\pandoc@box+\dp\pandoc@box\relax}%
  \Gscale@div\@tempb{\linewidth}{\wd\pandoc@box}%
  \ifdim\@tempb\p@<\@tempa\p@\let\@tempa\@tempb\fi% select the smaller of both
  \ifdim\@tempa\p@<\p@\scalebox{\@tempa}{\usebox\pandoc@box}%
  \else\usebox{\pandoc@box}%
  \fi%
}
% Set default figure placement to htbp
\def\fps@figure{htbp}
\makeatother





\setlength{\emergencystretch}{3em} % prevent overfull lines

\providecommand{\tightlist}{%
  \setlength{\itemsep}{0pt}\setlength{\parskip}{0pt}}



 
\usepackage[]{biblatex}
\addbibresource{../main/references.bib}


\usepackage{float}
\floatplacement{figure}{H}
\usepackage{booktabs}
\usepackage{colortbl}
\usepackage[table]{xcolor}
\definecolor{light-gray}{gray}{0.9}
\renewcommand{\arraystretch}{1.2}
\let\oldtabular\tabular
\let\oldendtabular\endtabular
\renewenvironment{tabular}{\small\oldtabular}{\oldendtabular}
\setlength{\tabcolsep}{4pt}
\makeatletter
\@ifpackageloaded{caption}{}{\usepackage{caption}}
\AtBeginDocument{%
\ifdefined\contentsname
  \renewcommand*\contentsname{Table of contents}
\else
  \newcommand\contentsname{Table of contents}
\fi
\ifdefined\listfigurename
  \renewcommand*\listfigurename{List of Figures}
\else
  \newcommand\listfigurename{List of Figures}
\fi
\ifdefined\listtablename
  \renewcommand*\listtablename{List of Tables}
\else
  \newcommand\listtablename{List of Tables}
\fi
\ifdefined\figurename
  \renewcommand*\figurename{Figure}
\else
  \newcommand\figurename{Figure}
\fi
\ifdefined\tablename
  \renewcommand*\tablename{Table}
\else
  \newcommand\tablename{Table}
\fi
}
\@ifpackageloaded{float}{}{\usepackage{float}}
\floatstyle{ruled}
\@ifundefined{c@chapter}{\newfloat{codelisting}{h}{lop}}{\newfloat{codelisting}{h}{lop}[chapter]}
\floatname{codelisting}{Listing}
\newcommand*\listoflistings{\listof{codelisting}{List of Listings}}
\makeatother
\makeatletter
\makeatother
\makeatletter
\@ifpackageloaded{caption}{}{\usepackage{caption}}
\@ifpackageloaded{subcaption}{}{\usepackage{subcaption}}
\makeatother
\usepackage{bookmark}
\IfFileExists{xurl.sty}{\usepackage{xurl}}{} % add URL line breaks if available
\urlstyle{same}
\hypersetup{
  pdftitle={Supplementary Legends for: Computational Tool Choice Impacts CRISPR Spacer-Protospacer Detection},
  pdfauthor={Uri Neri; Antonio Pedro Camargo; Rick Beeloo; Brian Bushnell; Simon Roux},
  colorlinks=true,
  linkcolor={blue},
  filecolor={Maroon},
  citecolor={Blue},
  urlcolor={Blue},
  pdfcreator={LaTeX via pandoc}}


\title{Supplementary Legends for: Computational Tool Choice Impacts
CRISPR Spacer-Protospacer Detection}
\author{Uri Neri \and Antonio Pedro Camargo \and Rick Beeloo \and Brian
Bushnell \and Simon Roux}
\date{}
\begin{document}
\maketitle


\section{Supplementary Materials}\label{supplementary-materials}

This legend-only version preserves the section structure, captions, and
explanatory notes from \texttt{supplementary.qmd} while omitting
embedded tables and figures. Source data remain available in
\texttt{draft/supp/figures/} and \texttt{draft/supp/figures/tables/}.

\subsection{Tables}\label{tables}

\subsubsection{Supplementary Table S1: Tool Configuration
Details}\label{tbl-supp-tool-config}

Parameters chosen for each tool to maximize recall on short spacers.
BLASTn settings emphasize ungapped alignments and permissive
identity/coverage to capture all valid matches; bowtie variants,
minimap2, MMseqs2, mummer4, sassy, strobealign, and x\_mapper entries
list versions and command-line templates. Full table is available in
\texttt{draft/supp/figures/tables/supp\_table\_S1\_tool\_config.tsv}.

\subsubsection{Supplementary Table S2: IMG/VR4 Contig Dataset Statistics
and Filtering Steps}\label{tbl-supp-contig-stats}

Summary of IMG/VR4 v1.1 high-confidence viral contigs after taxonomic
filtering and length cutoffs, including the final HQ benchmark subset
(fraction\_1). See
\texttt{draft/supp/figures/tables/supp\_table\_S2\_contig\_stats.tsv}
for the full table.

\subsubsection{Supplementary Table S3: CRISPR Spacer Dataset Composition
and Statistics}\label{tbl-supp-spacer-stats}

Statistics for the iPHoP spacer dataset before and after complexity
filtering, including length ranges, GC content, and removal criteria.
Data are provided in
\texttt{draft/supp/figures/tables/supp\_table\_S3\_spacer\_stats.tsv}.

\subsubsection{Supplementary Table S4: Detailed Statistics for All
Datasets}\label{tbl-supp-simulation-params}

\textbf{A. Real datasets} --- Counts and length/GC summaries for spacers
and contigs across IMG/VR4 fractions (0.0005 through 1).

\textbf{B. Simulated datasets} --- Equivalent summaries for all
synthetic datasets used in the study. Full details are in
\texttt{draft/supp/figures/tables/supp\_table\_S4\_real\_dataset\_sequence\_stats.tsv}
and \texttt{supp\_table\_S4\_simulated\_dataset\_sequence\_stats.tsv}.

\subsubsection{\texorpdfstring{Supplementary Table S5: Simulation and
job-generation configuration from
\texttt{Prepare\_all\_jobs.ipynb}}{Supplementary Table S5: Simulation and job-generation configuration from Prepare\_all\_jobs.ipynb}}\label{tbl-supp-prepare-jobs}

Per-dataset group configuration for simulations and tool runs, noting
search-space sizes, distance thresholds, and tool inclusion/exclusion.
See
\texttt{draft/supp/figures/tables/supp\_table\_S5\_prepare\_jobs.tsv}
for the structured table.

\subsubsection{Supplementary Table S6: Recall Values for Each Tool at
Different Mismatch Thresholds}\label{tbl-supp-recall}

Values are counted per spacer--contig alignment using exact distance for
each metric. \texttt{n\_found} tallies tool-reported positive alignments
at a given distance, \texttt{n\_total} is the ground-truth positive
count, and recall is \texttt{n\_found\ /\ n\_total}. Subsections report
IMG/VR4 and synthetic datasets separately. Complete values reside in
\texttt{draft/supp/figures/tables/supp\_table\_S6\_recall\_values.tsv}.

\subsection{Supplementary Figures}\label{supplementary-figures}

\subsubsection{Supplementary Figure S1: IMG/VR4 tool-vs-tool comparison
matrices}\label{fig-supp-real-matrix}

Pairwise tool-comparison matrices for IMG/VR4 fraction\_1 at exact
hamming distances (0--3). Each cell indicates regions detected by one
tool but not another; diagonals list per-tool totals. Original figure:
\texttt{draft/supp/figures/performance/imgvr4\_fraction1\_tool\_comparison\_matrix.svg}.

\subsubsection{Supplementary Figure S2: Simulated dataset recall vs
spacer occurrence frequency}\label{fig-supp-sim-recall}

High-occurrence-focused recall view for the high-insertion simulated
dataset (\texttt{ns\_500\_nc\_5000\_HIGH\_INSERTION\_RATE}) at hamming
distance ≤3. Original figure:
\texttt{draft/supp/figures/performance/imgvr4\_high\_ins\_recall\_vs\_occurrence.svg}.

\subsubsection{Supplementary Figure S3: IMG/VR4 fraction\_1 recall vs
occurrence frequency at hamming distance ≤3}\label{fig-supp-real-recall}

Recall versus spacer occurrence for the real IMG/VR4 fraction\_1 dataset
at hamming distance ≤3. Original figure:
\texttt{draft/supp/figures/performance/imgvr4\_fraction1\_recall\_vs\_occurrence.png}.

\subsubsection{Supplementary Figure S4: Comparison of simulated vs real
spacer characteristics}\label{fig-supp-spacer-comparison}

Distributions comparing simulated and real spacer features (k-mer
repeatability, entropy, base frequencies, GC, LCC). Panel A shows
filtered iPHoP spacers; Panel B shows fully simulated spacers
(ns\_100000\_nc\_20000). Original figures:
\texttt{figures/supp/filtered\_iphop\_spacers\_attribute\_distributions\_multiplot.png}
and \texttt{figures/supp/simulated\_spacer\_distributions.svg}.

\subsubsection{Supplementary Figure S5: Spacer characteristic
distributions by dataset}\label{fig-supp-spacer-distributions}

Spacer and contig characteristic distributions for real (IMG/VR4) and
simulated datasets, covering spacer size, mismatch counts, and
occurrence rates. Outliers were removed for fraction\_1 length
visualization. Original figure:
\texttt{figures/supp/real\_data\_contig\_distributions.svg}.

\subsubsection{Supplementary Figure S6: Distribution of matched spacers
from fraction\_1}\label{fig-supp-occurrence-distribution}

Distributions of matched spacers (hamming ≤3) across length, number of
occurrences, and per-distance alignment counts. Original figure:
\texttt{figures/supp/fraction\_1\_matched\_spacers\_distributions.png}.

\subsubsection{Supplementary Figure S7: Peak memory scaling (real and
simulated datasets)}\label{fig-supp-resource-memory}

Peak memory usage (RSS) scaling with dataset size. Panel A: IMG/VR4
subsamples; Panel B: simulated datasets. Axes are log-scaled; tool
identity encoded by color/shape. Original figure:
\texttt{figures/resource\_usage/resource\_usage\_memory\_2panel.svg}.

\subsubsection{Supplementary Figure S8: Tool recall consistency across
dataset sizes}\label{fig-supp-sim-consistency}

Per-tool recall consistency at hamming distance ≤3. Panel A: real
IMG/VR4 subsamples. Panel B: synthetic datasets. Panel C: heatmap of
per-tool recall across simulated runs. Original figures:
\texttt{figures/performance/imgvr4\_across\_fractions\_recall\_consistency\_hamming\_le3.svg},
\texttt{figures/performance/simulated\_cross\_fraction\_consistency.svg},
and
\texttt{figures/performance/simulated\_heatmap\_dataset\_consistency\_hamming\_le3.svg}.

\subsubsection{Supplementary Figure S9: Per-simulation non-planned match
counts (Hamming vs Edit)}\label{fig-supp-per-sim-background}

Grouped bars for each simulation (ns\_50000\_nc\_5000,
ns\_75000\_nc\_5000, ns\_100000\_nc\_10000, ns\_100000\_nc\_20000)
comparing hamming and edit distance non-planned match counts at
thresholds 0--5. Original figures:
\texttt{figures/ns\_50000\_nc\_5000\_background\_counts\_foldchange.svg},
\texttt{figures/ns\_75000\_nc\_5000\_background\_counts\_foldchange.svg},
\texttt{figures/ns\_100000\_nc\_10000\_background\_counts\_foldchange.svg},
and
\texttt{figures/ns\_100000\_nc\_20000\_background\_counts\_foldchange.svg}.

\subsection{Supplementary Notes}\label{supplementary-notes}

\subsubsection{Supplementary Note 1: Coordinate Tolerance Matching
Example}\label{note-supp-coord-tolerance}

Illustrative case from a synthetic dataset
(\texttt{ns\_100000\_nc\_10000}) where BLASTn, MMseqs2, and Sassy report
slightly different coordinates for the same spacer--contig locus. BLASTn
aligns positions 5003--5036 with 0 mismatches; MMseqs2 reports
5002--5036 with 1 mismatch; Sassy reports 5000--5036 with a 1 bp
insertion (CIGAR 34=1I1=); the planned occurrence is 5001--5036. All
represent acceptable matches at maximum hamming distance 1, motivating a
5 bp coordinate tolerance when aggregating alignments.

\subsubsection{Supplementary Note 2: Non-planned Match Rate
Analysis}\label{note-supp-nonplanned}

\paragraph{Hamming vs edit distance: per-simulation
results}\label{hamming-vs-edit-distance-per-simulation-results}

Per-simulation comparisons of non-planned match counts under hamming
versus edit distance for four fully synthetic datasets
(ns\_50000\_nc\_5000, ns\_75000\_nc\_5000, ns\_100000\_nc\_10000,
ns\_100000\_nc\_20000) at thresholds 0--5. Detailed counts are in
\texttt{distance\_metric\_analysis\_v3.ipynb} and associated CSV
exports.

\paragraph{Search-space-normalized non-planned match
rates}\label{search-space-normalized-non-planned-match-rates}

\subparagraph{Supplementary Table S7: Search-space-normalized metrics at
hamming/edit distance ≤3}\label{tbl-supp-normalized-threshold3}

Normalized non-planned match rates (per spacer per Gbp and per
Gbp\(^2\)) for each fully synthetic simulation. Edit distance ≤3
produces roughly 6--8× more non-planned matches than hamming ≤3. Full
numeric table:
\texttt{draft/supp/figures/tables/search\_space\_normalized\_metrics\_threshold3.csv}.

\subparagraph{Supplementary Table S8: Semi-synthetic non-planned match
rates by hamming threshold}\label{tbl-supp-semisynthetic}

Cumulative non-planned matches in the semi-synthetic dataset across
hamming thresholds (0--5), with rates normalized by spacer and contig
sequence space. Values are provided in
\texttt{draft/supp/figures/tables/supp\_table\_S8\_semi\_synthetic\_rates.csv}
and the corresponding TSV.

\paragraph{Caveats and interpretation}\label{caveats-and-interpretation}

Non-planned match estimates reflect compositional differences between
real and simulated sequences; semi-synthetic rates (real spacers ×
simulated contigs) exceed fully synthetic rates. For high thresholds
(≤4, ≤5) the counts are lower bounds because only heuristic tools were
feasible. These rates are most informative for relative comparisons
rather than absolute prediction.

\textbf{Complete analysis available in:}
\texttt{distance\_metric\_analysis\_v3.ipynb}, which includes
per-distance breakdowns, tool-specific validation, and detailed
comparison tables.

\subsection{References and Data
Availability}\label{references-and-data-availability}

Analysis notebooks, source code, and generated data are available in the
project repository: https://github.com/UriNeri/spacer\_matching\_bench.

Key notebooks: \texttt{spacer\_inspection.ipynb},
\texttt{distance\_metric\_analysis\_v3.ipynb},
\texttt{Prepare\_all\_jobs.ipynb},
\texttt{Performance\_simulated\_v2.ipynb},
\texttt{Performance\_imgvr4\_v3.ipynb}.

Datasets: filtered IMG/VR4 HQ contigs (fraction\_1 subset), filtered
iPHoP spacers, all synthetic datasets and ground truth files (Zenodo DOI
links in the main text).


\printbibliography



\end{document}
